\documentclass[10pt]{article}

% graphicx package, useful for including eps, pdf, and raster graphics
\usepackage{graphicx}
\DeclareGraphicsExtensions{.png,.png,.jpg,.pdf}

% basic packages
\usepackage{fancyhdr}
\usepackage{color}
\usepackage{parskip}
\usepackage{float}
\usepackage{microtype}
\usepackage{url}
\usepackage{hyperref}
\usepackage{booktabs}
\usepackage{makecell}
\usepackage{multirow}
\usepackage{pbox}
\usepackage{enumitem}
\usepackage[space]{xeCJK}

% fonts
\usepackage{fontspec}
\setmainfont{DroidSans}[
  Extension   = .ttf,
  UprightFont = *,
  BoldFont    = *-Bold ]
\setmonofont{DroidSansMono}[
  Extension   = .ttf,
  UprightFont = * ]
\setCJKmainfont{NotoSansCJK}[
  Extension   = .ttc,
  UprightFont = *-Regular,
  BoldFont    = *-Bold,
  Language    = Korean ]

\hypersetup{colorlinks,urlcolor=blue}
\urlstyle{same}

% page layout
\usepackage[top=1in, bottom=1in, left=0.75in, right=0.75in, includefoot, heightrounded]{geometry}

% Useful for seeing the layout as you tweak it.
%\usepackage{showframe}

\setlength{\headheight}{1.1in}
\pagestyle{fancy}
\fancyhf{}

\lhead{\includegraphics[height=1in,clip]{logo}}
\rhead{\pbox[b][1in][c]{\textwidth}{
    \small
    \textbf{다음 실험실이 실시한 검사:}\\
    Northwest Genomics Center\\
    University of Washington School of Medicine\\
    Genome Sciences Box 355065\\
    3720 15th Ave NE\\
    Seattle, WA 98195-5065}}

\newcommand{\PageLine}{\rule{\textwidth}{0.25mm}}

\newcommand{\link}[2]{\href{#1}{#2}\footnote{#1}}

\begin{document}

\begin{center}
\Large
\textbf{NWGC SARS-CoV-2 (COVID-19) Qual 폴리메라아제 연쇄 반응(Polymerase chain reaction, PCR) 실험실 보고서}
\end{center}

\bigskip

\begin{description}[font=\normalfont,align=left,labelwidth=6em]
\item [참여자 이름] \textbf{\VAR{pat_name|e}}
\item [생년월일] \textbf{\VAR{birth_date|e}}
\item [검체 식별자] \textbf{\VAR{qrcode|e}}
\item [샘플 제출 날짜] \textbf{\VAR{collect_ts|e}}
\item [결과 제공 날짜]
  %- if status_code == "never-tested"
  \textbf{테스트한 적이 없음}
  %- else
  \textbf{\VAR{result_ts|e}}
  %- endif
\item [검체 유형] 비강, 자가 채취
\end{description}

\PageLine

SARS-CoV-2 (COVID-19) Qual PCR 결과

%- if status_code == "negative"
\textbf{검출되지 않음}\\
%- elif status_code == "positive"
\textbf{SARS-CoV-2 검출(양성)}\\
%- elif status_code == "inconclusive"
\textbf{미확정}\\
%- elif status_code == "never-tested"
\textbf{검사 불가}\\
%- else
\VAR{0/0} %# poor man's raise
%- endif

\PageLine

참고:

\begin{itemize}
\item

  \textbf{본 프로젝트에 사용된 샘플 및 검사 유형은 King 카운티 및 기타 기관과 협력하여 공중 보건
  감시 프로그램의 일환으로 실시되었습니다.} 귀하가 받으신 검사는 의료진의 개인 의료 서비스를
  대체하지 않습니다.

\item

  \textbf{검체 관리 연속성:} 검사된 검체는 참여자가 자가 채취하여, 직접 라벨을 부착한 뒤 우편
  또는 배송 서비스로 SCAN에 제출되었습니다. 검체는 공인된 의료 전문가가 수집하고, 라벨을 부착한
  뒤 운송된 것이 아닙니다. 이는 검사 결과에 영향을 미칠 수 있으며 잘못된 음성 결과를 배제할 수
  없습니다. 추후 임상 검사가 이루어질 수 있습니다.

\item

  \textbf{검체 채취 장비:} 중비갑개(코안의 중간부분) 면봉 채취는 FDA/CDC의 긴급 사용
  승인(Emergency Use Authorization, EUA)을 받은 검체 유형이 아닙니다. 따라서, 검출되지
  않음/음성 결과가 나왔다고 해서 SARS-CoV-2 감염을 배제할 수 없습니다.

\item

  \textbf{``양성'' 결과는 Washington 주 보건부에 의해 확인될 수 있습니다.} 본 검사 결과에
  대해 질문이 있는 경우 담당 의료진과 상담해야 합니다.

\item{
  \textbf{귀하의 검사 결과가 ``검출되지 않음''인 경우 저희 실험실에서 SARS-CoV-2가 검출되지
  않았다는 뜻입니다.}

  어떤 검사도 완벽하지 않으며 음성 결과여도 코로나19(COVID-19)를 일으키는 SARS-CoV-2
  바이러스가 있을 약간의 가능성이 있습니다.  본 검사는 의료진이 검체를 수집하는 표준 방법을
  사용하지 않았으며, 샘플도 의료 서비스 환경에서 수집되어 운송된 것이 아니라 우편 또는 배송 서비스로
  제출되었습니다.  귀하의 검사 결과에 상관없이 몸이 좋지 않은 경우 집에 머물면서 다른 사람들과
  접촉하지 않는 것이 가장 좋습니다. 다른 건강 문제가 있거나 매우 아프거나 질병에 대해 걱정이
  된다면, 의료진에게 연락해야 합니다.  증상이 계속되는 경우, 코로나19(COVID-19)와 관련한 다른
  검사를 받는 것에 대해 의료진과 상담해야 합니다.
}

\item

  \textbf{검사 결과가 ``미확정''으로 나온 경우} SARS-CoV-2 표적에 대한 저희 실험실의 검사
  결과가 확실하지 않다는 뜻입니다.

  어떤 검사도 완벽하지 않으며 미확정 결과는 코로나19(COVID-19)를 일으키는 SARS-CoV-2
  바이러스가 있을 수도 있다는 의미입니다.  본 검사는 의료진이 검체를 수집하는 표준 방법을 사용하지
  않았으며, 샘플도 의료 서비스 환경에서 수집되어 운송된 것이 아니라 우편 또는 배송 서비스로
  제출되었습니다.  귀하의 검사 결과에 상관없이 몸이 좋지 않은 경우 집에 머물면서 다른 사람들과
  접촉하지 않는 것이 가장 좋습니다. 다른 건강 문제가 있거나 매우 아프거나 질병에 대해 걱정이
  된다면, 의료진에게 연락해야 합니다.  증상이 계속되는 경우, 코로나19(COVID-19)와 관련한 다른
  검사를 받는 것에 대해 의료진과 상담해야 합니다.

\item

  \textbf{검사 결과가 ``검사 불가''인 경우} 검체가 충분하지 않았거나 검사 과정에서 오류가
  있었다는 뜻입니다.

  검사를 시행할 수 없었기 때문에, 귀하에게 코로나19(COVID-19)를 일으키는 SARS-CoV-2 바이러스가
  있을 수도 있습니다.  본 검사는 의료진이 검체를 수집하는 표준 방법을 사용하지 않았으며, 샘플도
  의료 서비스 환경에서 수집되어 운송된 것이 아니라 우편 또는 배송 서비스로 제출되었습니다.  귀하의
  검사 결과에 상관없이 몸이 좋지 않은 경우 집에 머물면서 다른 사람들과 접촉하지 않는 것이 가장
  좋습니다. 다른 건강 문제가 있거나 매우 아프거나 질병에 대해 걱정이 된다면, 의료진에게 연락해야
  합니다.  증상이 계속되는 경우, 코로나19(COVID-19)와 관련한 다른 검사를 받는 것에 대해 의료진과
  상담해야 합니다.

\item

  \textbf{양성 결과(``검출'' 또는 ``바이러스 검출'')의 경우 검사 시설이 해당 주 및 카운티
  보건국에 결과를 보고할 것입니다.} 귀하는 이러한 결과를 직접 보고할 필요가 없습니다.  귀하의 검사
  결과가 신종 코로나바이러스에 양성으로 나온 경우 공중 보건 당국은 귀하의 질병 및 귀하가
  접촉했을지도 모르는 사람들에 대해 물어보기 위해 후속 조치를 취할 수 있습니다.

\item

  \textbf{이 검사 결과를 출력하거나 이메일로 담당 의료진에게 보내실 수 있습니다.}

\end{itemize}

\bigskip
\textbf{SARS-CoV-2 (COVID-19) Qual PCR 방법론:}

본 검사 분석은 역전사 PCR을 사용하여 SARS-CoV-2를 검출하도록 설계되었습니다. 이 분석은 체내
RNaseP 증폭을 통제하며 바이러스유전체 분절 Orf1B 및 S의 증폭을 사용하는 방법입니다.  이 검사는
University of Washington의 Northwest Genomics Center에서 개발하고 성능 특성을
정했습니다.  본 실험실이 개발한 검사는 검증을 받고 미국 Food and Drug Administration(FDA)에
긴급 사용 승인을 위해 제출되었습니다. 검사 타당화에 대한 FDA의 독립적인 검토는 아직 보류 중입니다.
본 실험실은 복잡한 임상 검사를 시행할 수 있도록 미국 실험실 표준 인증인 Clinical Laboratory
Improvement Amendments(CLIA)에 따라 인증되었습니다.

자세한 제공업체 정보는 Washington State Department of Health가 제공한 다음 링크를
참조하십시오.

\begin{itemize}
\item

  \link{https://www.doh.wa.gov/Portals/1/Documents/1600/coronavirus/COVIDcasepositiveKOREAN.pdf}{코로나19(COVID-19)
  확진 또는 의심 환자}

\item

  \link{https://www.doh.wa.gov/Portals/1/Documents/1600/coronavirus/COVIDexposed\_notSickYet\_KO.pdf}{코로나19(COVID-19)
  확진자에게 노출된 환자}

\item

  \link{https://www.doh.wa.gov/Portals/1/Documents/1600/coronavirus/COVIDsickNotExposedKorean.pdf}{코로나19(COVID-19)
  증상이 있으면서 확진자에 노출되지 않은 환자}

\end{itemize}

추가 자료는 Washington State Department of Health 및 CDC 웹사이트에서 확인할 수 있습니다:

\begin{itemize}
\item

  \url{https://www.doh.wa.gov/Emergencies/NovelCoronavirusOutbreak2020COVID19/Korean}

\item

  \url{https://korean.cdc.gov/coronavirus/2019-ncov/index.html}

\end{itemize}
\end{document}
